% 統計力学1 統計力学1-2 問題1（001-levels） 解いた日 2026-11-24
%% kotae
\toi{1}{○}{記述}
$(n_{x},n_{y},n_{z})=$

i) $(1,1,1)$
\[ E=E_{0}(1^{2}+1^{2}+1^{2})=3E_{0}\qquad\frac{E}{E_{0}}=3 \]
縮退度：1

ii) $(2,1,1)$，$(1,2,1)$，$(1,1,2)$
\[ E=E_{0}(2^{2}+1^{2}+1^{2})=6E_{0}\qquad\frac{E}{E_{0}}=6 \]
縮退度：3

iii) $(2,2,1)$，$(2,1,2)$，$(1,2,2)$
\[ E=E_{0}(2^{2}+2^{2}+1)=9E_{0}\qquad\frac{E}{E_{0}}=9 \]
縮退度：3

iv) $(3,1,1)$，$(1,3,1)$，$(1,1,3)$
\[ E=E_{0}(3^{2}+1^{2}+1^{2})=11E_{0}\qquad\frac{E}{E_{0}}=11 \]
縮退度：3
\naze{9 の次が 11 であること，すなわち和10 が $1,4,9$ から重複を許して3つ選んだ和では作れないこと（次の候補 $(2,2,2)$ は12）が式にも文にも無い（答案に必要な論点）。並べた4つが「低い方から4つ」である根拠を一言書く。値と縮退度は正しい。}
\toi{2}{◎}{}
\[ E_{0}=\qty{8.26e-38}{\joule} \]
\[ \Delta E=6E_{0}-3E_{0}=\qty{2.48e-37}{\joule} \]
\toi{3}{○}{記述}
\[ \frac{\Delta E}{\kB T}=\num{6.0e-17} \]
\[ \therefore\ \Delta E\ll\kB T \]
温度に対してエネルギー変化の間隔が十分小さいため，連続的に扱える。\nuke{$\delta U$ が選べる}
\naze{間隔が小さいことと「連続量として扱える」の間に，粗視化の条件 $\Delta E_{\mathrm{level}}\ll\delta U\ll U$ を満たす幅 $\delta U$ が選べる（幅 $\delta U$ の中に非常に多くの準位が入り，離散性が見えない）という一言が要る（答案に必要な論点）。数値は正しい。}
%% kizuki
\item 状態の数え上げでは，答えの準位だけでなく「その間に別の値が無い」ことまで示す。和の候補を $3,4,5,\dots$ と小さい順に当たり，作れない値（4，5，7，8，10）を飛ばしたことを一言書く。10 は $1+9$ に残り $0$ が，$4+4$ に残り $2$ が要るので作れない。
\item \mae{第1章 問題3 (4)・11-11} 「理由を述べよ」では，要点整理のどの条件に当てはまるかを名指しして書く。前回は $\Delta U=0$ の理由が抜け，今回は「間隔が小さい」から「連続的に扱える」へのつなぎが抜けた。ミクロカノニカル分配関数 $w$ は幅 $\delta U$ の中の固有状態の数で，$\Delta E_{\mathrm{level}}\ll\delta U$ なら幅の中に多数の準位が入り $w$ は滑らかになり，$\delta U\ll U$ なら $U$ はほぼ一点に決まる。この両方を同時に満たせることが「連続量として扱える」の中身である。
